\section{What Reliability Means in This Study}
A model's classification performance, its stability across a recording and the quality of its probability scores describe different properties. We evaluate balanced accuracy and AUROC against reference labels, label changes across defined windows or recording states, unanimous correct and wrong classifications, and the dependence of decisions on score thresholds. None alone establishes clinical reliability. A constant-label classifier can show perfect stability yet have 50\% balanced accuracy when both binary classes are present. The three main studies below share this concern, but their stability definitions and sampling units differ.

\section{Study A: Temporal Decisions and Persistent Errors}
The original NEMAR Study A used a class-balanced spectral SVM and compared five consecutive 60-second windows. The first-minute participant-level balanced accuracy was 73.48\% for AD versus CN (46 participants) and 73.72\% for FTD versus CN (41 participants). Across five repeated evaluations, the AD/CN task yielded 108/230 stable-correct, 20/230 stable-wrong and 102/230 changed-label participant-repeat outcomes. The FTD/CN task yielded 109/205 stable-correct, 27/205 stable-wrong and 69/205 changed-label outcomes. Consequently, the label-change fractions were 44.35\% and 33.66\%, while stable-wrong fractions were 8.70\% and 13.17\%. These correlated repeat rows cannot be treated as 230 or 205 independent people.

\ThesisFigure[Five-minute temporal accuracy and prediction states]{fig_07_01_temporal.pdf}{Historical NEMAR Study A SVM results. First-minute balanced accuracy uses 46 AD/CN and 41 FTD/CN participants. The stable-correct, stable-wrong and changed-label totals use correlated participant-repeat states (230 and 205). The repeated folds do not create independent patients and the results do not measure DSRF temporal performance.}{fig:temporal}

Figure~\ref{fig:temporal} demonstrates that a reasonable first-minute group score can coexist with substantial time-related disagreement and persistent mistakes. It does not show that longer recordings necessarily improve accuracy or that EEGraft fixes the original SVM's temporal variation. The independent historical duration audit was partly reconciled as a model-counting error (690 rows per model, 1,380 combined); a complete corrected downstream rerun remains outside the submitted evidence. The archived duration-position plot is therefore omitted from the main thesis, and no optimal recording duration is asserted.

\section{Independent Isfahan MCI Negative Control}
The original Isfahan MCI-versus-healthy SVM experiment screened 61 recordings and retained 49 after a documented quality gate. Release EEG\_1 retained 25 and EEG\_2 retained 24, with separate nested validation. First-minute balanced accuracy was 40.78\% and 55.83\%, respectively. Nonetheless, EEG\_1 labels were stable across the five-minute comparison for 94.40\% of its correlated repeat records, while 59.20\% were stable but wrong. EEG\_2 stability was 80.00\%, with 35.83\% stable but wrong; change rates were 5.60\% and 20.00\%. This is a negative control about the meaning of consistency, not evidence of an external EEGraft model prediction for MCI.

\ThesisFigure[Isfahan MCI stability negative control]{fig_07_02_mci.pdf}{Independent release-specific Isfahan MCI-versus-healthy SVM experiment (61 recordings screened; 49 retained). Reported first-minute balanced accuracies of 40.78\% and 55.83\% coexist with high stable-but-wrong fractions. The 125 and 120 state records are correlated across folds, and the two releases are not pooled or interpreted as EEGraft transfer.}{fig:mci}

A small flip rate is useful only when it is accompanied by correct decisions. This example is important because a naive strategy that rewards stability without checking labels could prefer an almost constant but inaccurate classifier.

\section{Study C: Whole-Recording Unanimity on the Locked Test}
The 57 CAUEEG binary test recordings each contributed 24 sampled windows for the reported unanimity analysis at the original threshold of 0.50. For DSRF, four recordings had 24 unanimous correct predictions, zero had 24 unanimous wrong predictions and 53 were not unanimous. For XGBoost, 13 were unanimous correct, 12 were unanimous wrong and 32 were not unanimous. Thus XGBoost looked more consistent under this strict all-window rule while also showing twelve completely stable errors. DSRF was less often unanimous, but the absence of unanimous wrong records does not by itself make it the more reliable classifier. DSRF correctly classified 39/57 recordings using mean-score aggregation, whereas XGBoost classified 28/57 correctly under the original rule.

\ThesisFigure[Unanimous correct and wrong external predictions]{fig_07_04_unanimous.pdf}{CAUEEG locked binary test (57 recordings), original 0.50 thresholds, 24 sampled windows per recording. DSRF: 4 unanimous correct, 0 unanimous wrong, 53 not unanimous; XGBoost: 13 unanimous correct, 12 unanimous wrong, 32 not unanimous. ``Unanimous'' is an all-window test and differs from an eyes-open/closed condition-pair flip.}{fig:unanimous}

\section{Threshold Choice Changes Apparent Consistency}
The post-hoc CAUEEG event audit reconstructs mean model scores separately over eyes-open and eyes-closed windows within each eligible development recording. At least two windows are required for each state. There were 404 eligible training recordings and 42 eligible validation recordings, not the entire 950 and 82 partitions. At the original 0.50 threshold, 142/404 (35.15\%) DSRF training recordings changed their predicted label between condition means. XGBoost changed for 29/404 (7.18\%). However, with the XGBoost training-selected threshold at 0.70, its change count became 134/404 (33.17\%); in validation it increased from 5/42 (11.90\%) to 15/42 (35.71\%). DSRF's validation change rate at 0.50 was 13/42 (30.95\%). The score sequences and model weights were unchanged during these comparisons.

\ThesisFigure[CAUEEG threshold sensitivity and calibration]{fig_07_03_threshold_results.pdf}{Frozen CAUEEG model results. Training and validation panels compare labels assigned to condition-averaged eyes-open and eyes-closed scores at original versus training-selected thresholds. Locked-test balanced accuracy and Brier error are reported separately. The 0.70 XGBoost rule and Platt calibration are post-hoc exploratory variants, not model retraining or a repeated independent test.}{fig:threshold-results}

Figure~\ref{fig:threshold-results} is central to the thesis title: threshold-aware assessment is needed because apparent consistency can move by roughly 26 percentage points for the same XGBoost scores. The original DSRF-minus-XGBoost training flip-rate difference was 27.97 percentage points; when XGBoost uses 0.70, it is only 1.98 points. Likewise, calibration changed model-specific flip rates, reducing any simple conclusion about architectural stability. These outcomes were examined after the original locked external results and are exploratory.

\section{Recorded Eye State, Timing and Score Differences}
The upstream event audit parsed 46,260 event entries and aligned saved model-window scores with event labels in 593/950 training and 58/82 validation recordings under a conservative eligibility rule. Of these, 404 training and 42 validation recordings had at least two windows in both eye states. In the eligible training sample, mean eyes-open minus eyes-closed score shifts were +0.1038 for DSRF and +0.0533 for XGBoost. Validation shifts were +0.0787 and +0.0398. The median open-minus-closed time offsets were approximately $-95.1$ seconds in training and $-125.9$ seconds in validation, so eyes-open and eyes-closed intervals were not balanced in time.

A later within-recording linear time adjustment estimated exploratory training eye-open coefficients of +0.1254 for DSRF and +0.0650 for XGBoost; their reported exploratory training intervals excluded zero. Validation coefficients were +0.0697 and +0.0469, but their 95\% intervals included zero. These findings suggest score differences are associated with recorded conditions, but they \emph{cannot establish a causal eye-state effect}: time ordering, event selection and other acquisition factors remain possible explanations. They should not be used to prescribe clinical eye-state protocols.

\ThesisFigure[Eye-state and elapsed-time audit]{fig_07_05_eye_audit.pdf}{Exploratory CAUEEG development-only eye-state audit based on 404 training and 42 validation recordings with sufficient annotated eyes-open and eyes-closed windows. Score shifts and time adjustment are within-recording observational associations. The original locked test waveforms, events and predictions were not accessed during Codes 21--23.}{fig:eye-audit}

\section{Age Differences and Incremental Score Information}
The CAUEEG development data show a strong demographic imbalance. In training, Normal recordings had mean age 65.37 years and Dementia recordings 76.59 years, a difference of 11.23 years; in validation the corresponding means were 65.97 and 79.40, a difference of 13.43 years. A model using only age achieved validation AUROC 0.8397 and balanced accuracy 74.26\%. Adding frozen DSRF scores to linear age increased validation AUROC to 0.8515 ($+0.0118$), while adding XGBoost increased it to 0.8471 ($+0.0074$). A later recording bootstrap gave an exploratory 95\% interval of $-0.0860$ to $+0.1081$ for the DSRF incremental validation AUROC. It spans zero. The corresponding flexible-age comparison also had a validation interval crossing zero.

\begin{table}[!htbp]\centering\small
\caption{Exploratory CAUEEG development validation: age-only versus age-plus-frozen-score classifiers}\label{tab:age-validation}
\begin{tabular}{@{}lrrr@{}}\toprule
\textbf{Validation model}&\textbf{AUROC}&\textbf{BAcc (\%)}&\textbf{Binary Brier}\\\midrule
Age only&0.8397&74.26&0.1659\\
Age + DSRF score&0.8515&77.79&0.1613\\
Age + XGBoost score&0.8471&78.24&0.1555\\\bottomrule
\end{tabular}\end{table}

In training cross-validation, adding DSRF scores to linear age gave an apparent AUROC gain of +0.0457, with a conditional recording-bootstrap interval from +0.0251 to +0.0677. However, such a development improvement is not a fresh independent test. The small validation increment and its interval crossing zero are more appropriate constraints on an incremental-value claim.

\ThesisFigure[CAUEEG age confounding and incremental EEG value]{fig_07_06_age_audit.pdf}{Post-hoc CAUEEG training and validation age audit. Dementia-labeled records are older than Normal records by 11.23 years in training and 13.43 years in validation. The age-only validation AUROC is 0.8397; the additional frozen DSRF score changes it by +0.0118, with a recording-bootstrap 95\% interval of $[-0.0860,0.1081]$. The interval crosses zero; incremental EEG information beyond age is not established.}{fig:age-audit}

This audit does not imply that the EEG decoder uses age as a direct input. It indicates a confounding route: group labels and age are associated, so external EEG performance cannot be interpreted as age-independent discrimination without additional controls. The analysis is not a substitute for age-matched prospective validation. Its unit is the recording and its confidence intervals do not guarantee unique-patient sampling.

\section{Study B: Artificial Signal Perturbations and Electrode Sensitivity}
The original 52-person development stress tests retained DSRF BAcc 77.06\% on clean input and reported mean BAcc of 71.87\% at 10\% Gaussian noise, 70.07\% at 20\% noise, and 76.78\% when a single EEG channel was zeroed. The corresponding average label-change rates were 23.08\%, 27.69\% and 7.31\%. Five random seeds describe repeated perturbation runs, not five new subject groups. These values are already shown in the combined ablation and robustness Figure~\ref{fig:ablations}; no duplicate bar charts are added here.

A separate electrode-occlusion audit replaced a channel within each window by its within-window mean and measured the signed change in true-class score. For the 52 out-of-fold participants, per-class ranks differ; the complete channels and signs are reported in Appendix Figure~\ref{fig:occlusion}. Input masking can introduce distribution shifts and does not identify causal neural pathology or validated disease biomarkers. The technical stress experiments are distinct from natural recording-event sensitivity in CAUEEG.

\section{Reliability Synthesis}
Together, the reliability experiments provide a concrete reporting result. First, pair consistency statistics with their correct and incorrect counts, because a stable label can still be a stable error. Second, retain a model's original threshold when reporting primary results and show how an independently chosen secondary rule changes decisions without changing model weights. Third, report calibrated probability quality separately from accuracy and unanimity. Finally, record the evaluation unit, eye-state timing and age distribution when comparing cohorts. These practices follow directly from the observed NEMAR, CAUEEG and Isfahan results and can be reused in future EEG model evaluations; they do not require treating the different models or studies as one pooled experiment.
